
%%%%%%%%%%%%%%%%%%%%%%%%%%%%%%%% SCRIPT FOR E-RARA AND MDZ FILES     %%%%%%%%%%%%%%%%%%%%%%%%%%%%%%%%%%%%%%%%%%%%%%%%
%%%%%%%%%%%%%%%%%%%%%%%%%%% fini le 30.04.2025 par F. GOY            %%%%%%%%%%%%%%%%%%%%%%%%%%%%%%%%%%%%%%%%%%%%%%%%
% !TeX TS-program = lualatex
\documentclass{article}
\usepackage[T1]{fontenc}
\usepackage{microtype}
\usepackage[pdfusetitle,hidelinks]{hyperref}

\usepackage{polyglossia}
\setmainlanguage{english}
\setotherlanguages{latin,greek}
\usepackage[series={},nocritical,noend,noeledsec,nofamiliar,noledgroup]{reledmac}
\usepackage{reledpar}

\usepackage{fontspec}
\setmainfont{TeX Gyre Termes}

\usepackage{sectsty}
\usepackage{xcolor}

\usepackage{fancyhdr}
\pagestyle{fancy}
\fancyhf{}
\fancyhead[LE,RO]{\nouppercase{\leftmark}}  
\cfoot{\thepage}
\renewcommand{\headrulewidth}{0.4pt}

% Redefine \section to remove numbering
\usepackage{titlesec}
\titleformat{\section}[block]{\normalfont\scshape\color{gray}}{}{0pt}{} % no number in heading
\titleformat{\subsection}[hang]{\normalfont}{}{0pt}{} % also remove subsection number
\titleformat{\subsubsection}[hang]{\normalfont\footnotesize\color{black}}{}{0pt}{}

% Modify how section marks are stored to exclude numbers
\makeatletter
\renewcommand{\sectionmark}[1]{%
	\markboth{#1}{}} % Only store the section title, without number
\renewcommand{\subsectionmark}[1]{%
	\markright{#1}} % Only store the subsection title, without number
\renewcommand{\numberline}[1]{} % Hide the section number in TOC
\makeatother

\begin{document}

\date{}
        \title{In priorem epistolam Pauli ad Timotheum commemtarii : [Thomas Cajetan], [1531]}
\maketitle
\tableofcontents
\clearpage
\begin{pages} 
\beginnumbering
        
\section*{PRIORIS }
\section*{AD TIMOTHEVM }\pstart missis inuenitur huiusmodi integra salutatio in fine nisi in epistola ad colo. in qua legitar. gratia cum omnibus vobis amen, et non dicitur domini nostri iesu christi etc. ¶ Et apud graecos inuenit subscri- ptio talis. Ad thessalonicenses secunda scripta athenis. ¶Caietae die. XVII. martii. M.D.XXIX.  \pend
\phantomsection
\addcontentsline{toc}{subsection}{\textit{THOMAE DE VIO CAIETA- NICARDINALIS SANCTI XISTI. IN PRIOREM EPISTOLAM PAVLI AD TI- MOTHEVMCOMMENTARII. }}
\subsection*{\textit{THOMAE DE VIO CAIETA- NICARDINALIS SANCTI XISTI. IN PRIOREM EPISTOLAM PAVLI AD TI- MOTHEVMCOMMENTARII. }}\pstart \huge\textbf{P}\normalsize AVLVS apostolus iesu christi sum imperium del saluatoris nostri et Deest. domini iesu christi spei nostrae.] Non solum se apostolum iesu chri- sti explicat sed apostolatum ipsum explicat non quaesitum sed iniunctum sibi a deo saluatore nostro: naturam exprimens diuinam et officium sal- uandi nos. vt intelligamus iniunctum ei fuisse apostolatum ad exequen- dum officium salutis nostrae. Nec tantum a deo sed etiam ab homine do- mino iesu christo iniunctum sibi explicat. Et dicit domini, eo quod ad statum dominii iam peruenerat iesus christus quum iniunxit apostolatum pau- lo. Adiungit que  spei nostrae, eo quod christus imortalis iam et gloriosus sum corpus  est spes nostrae futurae bearitudinis. christi nanque  resurrectioni in- nititur spes nostra quod erimus quoque  nos immortales sum corpus. TIMO- THEO dilecto. pro. germano.] Eadem dictio haberur hic quae habita est ad philip. dicendo rogo te germane (seu germana )compar. Quadrat autem dictio illa tum aequali tum minori. vnde et ibi dictum est germana compar  et hic dicitur ger- mano filio. Significat enim vere ac confone esse talem. vnde germanus filius est qui vere est consonans patri filius. FILIO in fide.P ad differentiam filii naturalis. vt enim patet act. 1o.paulus asciuit sibi timo- theum. quem ex his verbis apparet a paulo fuifle conuersum ad fidem:alioquin non appellaret ipsum fi- lium in fide. GRATIA et misericordia et pax.] Superfluunt coniunctiones. Noua salutatio. nam ha- ctenus salutauerat gratia et pax  modo tria dicit gratia misericordia pax, ea ratione quod praelatis maxime necessaria est misericordia:tanq̃ loco dei cui proprium est misereri  praesint. A DEO patre et Deest. domi- no  iesu christo domino noltro. ] Bis dicit domino. Iemel ablolute  domino iesu christoriterum relati- ue ad nos domino nostro. Est siquidem absolute dominus omnium, et est peculiariter dominus noster tanq̃ sponte seruorum.  \pend\pstart \huge\textbf{S}\normalsize ICVT rogaui te. ] Completa salutatione inchoat tractatum epistolarem. Vbi aduerte quod aduer- bium sicut non habet in textu aliquid correspondens ad complendam orationem vsque  ad calcem huius capituli, vbi dicitur  hoc praeceptum commendo tibi. ita quod illic compler paulus orationem quam hic inchoat. Et est sensus. sicut rogaui te etc.  ita hoc praeceptum commendo tibi  quemadmodum olim rogaui te nunc tibi commendo hoc praeceptum. Vnde omnia interiecta ab hinc illucusque   vere inter- iecta sunt ad explicandum necessitatem praecaepti  et quod olim dederat et quod modo commendat. VT remaneres ephesi quum irem in macedoniam, vt denunciares. pro. vt praeciperes  quibusdam ne aliter docerent.] Dictio graeca interpretata aliter docerent composita est ex alio et doctor : varietatem doctorum clare significans. Mandauit paulus timotheo vt praeciperet quibusdam ne essent varii do- ctores: quod esset variam inducere doctrinam. Et quantum subiuncta insinnant, aduersus illos loqui- tur qui aliam de legalibus doctrinam inducebant. nam tales doctores explicabit paulo infenus. NE- QVE intenderent fabulis.] Thalmudicas fictiones prohibet. ET genealogiis interminatis.] Non dicit neque  intenderent genealogiis  sed addit interminatis : quae scilicet sunt sine termino in sacris sen au- tenticis scripturis, multae enim genealogiae tunc temporis erant interminatae. Et rationem huius vltimi subdit. QVÆ quaestiones praestant magis quam aedificationem. pro. dispensationem  dei quae est in fide] Interminatae siquidem genealogiae quod praebeant quaestiones  clarum est. Et q magis multiplicandis quae- stionibus prosint quam diuinae dispensationi seu gubernationi quam exercet deus in fide  paulus docet. nam ad dei gubernationem qua gubernat ea quae sunt fidei  parum aut nihil prodest notitia huiusmodi genealogiarum : vt patet de multis genealogiis quas hemmus in libris paralipo. ⁋ FINIS autem precaepti est di lectio. ] Dixerat vt preciperes quibusdam ne vati essent doctores, perperea vnitatem finis omnium precaeptorum legis  \pend
\section*{CAPVT }
\section*{.I. }
\marginpar{[ p.140 ]}\pstart explicat:vt ex ipsa vnitate finis praecaeptorum  varietas arguatur doctorum. Finis inquit praecaepti est dilectio, non qualiscunque  sed procedens a triplici fonte. Primus est. DE CORDE puro.] hoc est affectu puro  a passionibus et a malicia. ET conscientia bona.] Secundus fons est applicatio scientiae  applicatio inquam bona, ad agendum  ad cauendum etc. Multi n.hñt applicationem sciae malam: et innumera mala ĩde se- quunt. ET fide non ficta.] Tertius est hic fons vltimo loco scriptus : eo quod duo primi naturae sunt, hic autem tertius reuelarae est gratiae. Et dicit fide non ficta,  propter hypocritas in fide: qui fingebant se christianos  et nec erant christiani nec iudaei.  ¶ Discurre per praecaepta : et inuenies omnia ordinari ad dilectionem talem. huiusmodi enim dilectio est dilectio dei et proximi ex vera fide in intellectu  et ex puro corde in affectu  et ex conscientia bona in applicatione et exequutione : quae perficit et mandata primae et se- cunde tabulae ⁋Haec sunt quae paulus proponit tractanda in hac epistola. Verum antequam sigillatim par- riculas propositas tractet, necessirarem commemorat remansionis timothei ephesi ad precipiendum quae dixit ex casu aberrantium ab his. Vnde dicit. A QVIBVS.] praecaeptis et conditionibus finis praecaeptorum. QVIDAM aberrantes. ] Similes pseudapsis significat. CONVERSI sunt in vaniloquium. ] ab vtili do- ctrina. VOLENTES esse legisdoctores. ] Hinc apparet iudaeos esse istos qui doctrinam legis moysi i- gerebant. Et hinc coniicitur istos esse quos appellauit alia docentes. Et iungitur vaniloquio non do- ctrina legis, sed voluntas vtendi officio doctoris legis versa dicitur in vaniloquium. Et ratio subditur. NON intelligentes neque  quae loqurntur.] ipsimet praecaepta legis. NEQVE de quibus affirmant.] obligari ad praecaepta legis.  \pend\pstart ¶ SCIMVS autem quod bona est lex.] Intendit paulus commemorare vt timotheus non solum praecipiat vi- tari fabulas et quaestiones de genealogiis  sed ent doctores legis moysi male applicantes illam: quod esset si ap- plicarent illam ad obligandum iustificatos in chr̃o. Et propter hoc non repraehendit legem sed laudat eam : dicendo scimus autem quod lex bona est. Vituperat autem male applicantes eam subiumgendo. SI QVIS ea legi- time vtat̃. P hoc est si ꝗs eam sum ipsammet legem applicet. Et ne vageris quaerendo quod  est legitime appli- care  explicat subiungendo. SCIENTES. pro. sciens  hoc.] Ecce legitimus vsus  sciens hoc. QVIA. pro. quod lex iusto non est posita ] Formalis est sermo pauli.nam quum lex non superflue ponat̃  et iustus quatenus iustus non indigeat lege, formaliter venficatur quod iusto non est lex scripta posita. Rursus quum medium le- gis sit poena et iustus quatenus iustus nulli sit obnoxius poenae, formaliter verificat̃ quod iusto non est lex po- fita. ⁋Et intendit per hoc quod quemadmodum male applicatur lex ad iustos, ita male applicat̃ ad iustificatos in chr̃o. et propterea subiunget de seipso iustificato in christo. Cum qua veritate stat quod si iustus recedat a iu- stiriaistarim inuenit in ordine eorum quibus est lex posita: vtpote non iustus. SED iniustis.] E regione iusti vniuersaliter  nomniat per partes multa genera iniustor. Et primo iniustos in genere. ET non subditis. pro. et iis qui subdi nolunt, ipiis.] in patriam ac parentes etc.  ET peccatoribus. ] a recto recedentibus. SCELE- RATIS. ] qui prae aliis copia flagitiorum exuberant. ET contaminatis. pro.prophanis.] ea q̃ religionis sunt polluentibus. PATRICIDIS et matricidis.) Clara sunt crimina haec. HOMICIDIS  fornicariis! masculorum concubitoribus, plagiariis.] Dictio plagiariis significat fures hoium in mancipia, vel ven- denda vel redimenda ab ipsismet captis sen sibi coniunctis. MENDACIBVS et periuris. Superfluir et. et si quid aliud sanae doctrinae aduersat̃. ] Ne per omnes species discurrat  vnico verbo reliqua compraehendit. QVÆ est. ] Superflunnt hae duae dictiones adiunctae ab aliquo qui restringere voluit sequentem particu lam ad sanam doctrinam:quum tamen possit alio referri.nam et pont referri ad id quod dixerat  sciens hoc : vt mon- explicet sciens hoc. SECVNDVM euangelium gloriae beati dei.]  ad significandum vnde scit̃ haec libertas iusti a lege. s. exeuangelio gloriae beari dei. et ad discernendum vere scientes ab illis qui volunt esse legis- doctores et non intelligunt quod  dicant, quam nesciunt euangelium gloriae dei beati. ¶Potest nihilominus re- terri ad sanam doctrinam, ad explicandum quae est sana doctrina: quae scilicet est sum euangelium gloriae non hominis sed dei  non sum statum praesentem sed beatum. Et apte loquens de lege  appellat euange- lium euangelium gloriae beati dei, quoniam finis iusti qui vere iustus est  est gloria apud deum beatum quam oportet esse gloriam foeliciratis aeternae. QVOD creditum est mihi. ] Ad declarandum exfacto in seipso pau- lo quod dixerat (videlicet lex iusto non est posita) subiungit qualis modo est in statu euangelii et qualis prius fuerit  et vnde etc.  Ab officio itaque  euangelico inchoat: et vt materiam gratiarum actionis suppu- tat multitudinem beneficiorum sibi collatorum ad enangelicum munus  subiungendo. ¶ GRATIAS ago. pro. et gratiam habeo  ei qui me confortauit. pro. potentem fecit  in christo iesu domino nostro.] Superfluit in. Et gratiam habeo ei domino nostro iesu christo. vel gratiam habeo ei  qui fecit me poten- tem christo iesu domino nostro. non fecit me potentem mundo aut mihi ipsi sed christo: hoc est ad serui- tium christi. Ꝙ FIDELEM me existimauit. Ecce aliud beneficium quod me iudicauit fidelem  quod me vsus est vt fideli. PONENS in ministerio.] procul dubio suo ad propagationem regni coelorum. QVI prius blasphemus fui ] Explicatis donis in libertate euangelii, explicat crimina quoque  propria in sta- tu legis. Et dicit se fuisse blasphemum  propter iniurias quas irrogauerat christo. Agnoscens siquidem chr̃m esse verum deum  agnoscit peccatum suum in chr̃m fuisse blasphemiam. ET psequutor.] subaudi ec- clesiae christi. ET contumeliosus. pro. et violentus] Non solum persequutus fuerat paulus eccelesiam  \pend
\section*{PRIORIS }
\section*{AD TINOTHEVM }\pstart christi, fed et violentiam intulerat: vt paret in actibus apostolorum. SED. Deficit et misericordiam dei con- sequutus sum.] Superfluir dei. QVIA ignorans feci in incredulitate.] Romnem ex parte qualitatis pec- cati reddit quare misericordiam consequutus sit  quia ignoranter peccauit in statu incredulitatis. Ita quod ignorans refert̃ ad vtrunque  : tum ad feci tum ad ĩcredulitatem. nam et peccatum ĩcredulitaris ex ignorantia in paulo fuit : et factum blasphemiae  persequutionis et violentiae ex ignorantia quoque  fuit. In cuius signum quam primum audiuit ego sum iesus nasarenus quem tu persequeris  statim dixit domine quid me vis facere: Nonnulla siquidem ratio veniae est ex ignorantia peccare: quoniam affectus peccantis ex ignorantia  non tendit per se ad malum: vtpote qui non tenderet in illud si sciret esse malum. SVPERABVNDA- VIT autem gratia domini nostri.] Non solum misericordiam consequutus sum condonantem mihi peccata mea, sed etiam superabundauitgratia domini nostri: vtpore promouens me ad apostolatum. CVM fide et dilectione quae est in chro iesu. ] Intellige superabundauit  referri non solum adgratiam, sed etiam ad do- na fidei et dilectionis in chr̃o. haec .n. tria abundanter in paulo apparent. nam et gratia apłica in eo exu- berat: et fides valde inconcussa tum in perpessis tum in miraculis apparet:dilectio quoq ranta vt ipse dicat. quis infirmatur et ego non infirmor  quis scandalizatur et ego non vror etc.   \pend\pstart FIDELIS sermo et omni accaeptione dignus.] Ne libertatem iusti ex seipso tantam ex quodam singu- lari astruere videatur, subiungit vniuersalem causam huiusmodi libertatis  aduentum chr̃i. Et sermonem huius causae communis statui euangelico praemitrit esse tum fidelem  hoc est veracem, tum omni accaeptatio- ne dignum: vt tam parti cognoscitinae congruere significetur ex veracitate  quam parti affectiuae ex accaeptabi- litate. QVIA. pro. quod  christus iesus venit in hunc mundum peccatores faluos facere.] Superfluit hunc. Ecce sermo fidelis et omni accaeptione dignus. Et dicendo venit in mundum  diuinam iesu chrsti perso- nam significat. Quum enim esset supra mundum, ingressus est mundum  factus est vna mundi pars quum factus est homo.¶ Et dicendo peccatores saluos facere  explicat quorsum venerit. E regione ad legem (quae posita est ad puniendos peccatores) dicitur quod iesus christus venit ad saluandum peccatores. quod est eri- pere eos a subiectione legis. QVORVM primus ego sum. ] Relatiuum quorum  non refert peccatores absolute, sed peccatores saluatos a chr̃o. Et est sensus. quorum peccatorum qui saluantur per christi aduen- tum ego sum primus : hoc est precipuus. pctonrum siquidem tunc temporis saluatorum per christum  duo fuerunt ge- nera. Alterum peccantium ex infirmitate: vt mulier q̃ erat in ciuitate peccatrix  vt petrus qui ex timore ne- gauit et similes. Alterum eorum qui per ignorantiam crucifixerunt chr̃m  qui postea adueniente spiritu sancto poenitentiam egerunt : vt dicitur in actibus apostolorum. Paulus his omnibus saluatis peccatoribus prestat, tum quam et blasphemus et psequutor et violentus fuerat. tum quia nec per aduentum spiritus scti et predicationem apostolorum conuersus est vt alii. Et propterea merito inter saluatos peccatores primum (hoc est precipuum) se dicit. ET ideo misericordiam consequutus sum: vt in me primo ostenderet iesus chr̃s oemm patientiam pro. longanimitatem, ad informationem eorum qui credituri sunt illi. pro. in illum  in vitam aeternam. ] Primo non est aduerbium sed nomem. et significat precipuum. Reddiderat prius romnem consequutae prius mise- ricordiae ex parte qualitatis peccati sui : modo aliam romnem reddit ex parte exempli tribuendi aliis pec- catoribus. Ideo inquit misericordiam consecutus sum vt in me primo (hoc est precipuo) monstraret iesus chr̃s oemm longanimitatem propriam quam seruat expectando peccatorum pnĩam (animi .n. in longum propensi est hmoni expectatio) ad statuendum me exemplar aliis qui credituri sunt. Informationem intellige exem- plari forma. Et est sensus. vt in me praecipuo ostenderet longanimitatem ad exemplarem informationem itineris in vitam aeternam respectu eorum qui credituri sunt in chr̃m. In ipreo siquod em paulo habemus prae ocu- lis exemplar eundi in vitam aeternam  quamuis peccatores simus. Modus autem eundi explicatur per fidem in chri- stum. Et haec est propria ratio propter quam paulus precipuus inter peccatores saluatos misericordiam consecutus est. ¶ REGI autem saeculorum. ] Superius dixerat   etgratiam habeo : mon exercet gratiarum actiones, dicendo. regi autem non vnius saeculi sed saeculorum. IMMORTALI. ] ad differentiam regum longaeuae vitae : puta noe qui et saecu- lo añ diluuium prefuit et saeculo post diluuium, qui tñ non fuit immortalis. INVISIBILI.] Hoc ad differen- tiam coelestium corporum dicit̃ coelum. n. licet regat saecula et sit immortale, est tum visibile. SOLI deo.] Defi- cit sapienti. legendum est. SOLI sapienti. ] suapte natura. reliquae .n. substantiae immortales ac inuisibiles quum in nudis naturalibus desipere possint peccando, non sunt suapte natura sapientes: sed quaecunque  sunt sapien- tes  ex liberali dono dei sapiunt. Quum dixisset tot regis conditiones, par fuit vt plenae sapiae meminerit : ex- plicando non solum quod est sapiens  sed quod est solus sapiens. DEO. ] Faciendum est punctum post soli sapienti, et deinde subiungendum deo. quadrat siquidem post tot adiectiua vt natura ipsius regis explicet̃, dicendo deo. HONOR et głia in saecula saeculorum amen. ] Haec de ĩteriectis. Vnde et rediens ad principium epistolae sicut rogaui te vt remaneres ephesi vt praeciperes quibusdam etc.  subiungit. HOC precaeptum commendo tibi fili timothee.] Illo mon quo rogani te vt remaneres et preciperes ephesi : nunc quoque  commendo tibi hoc precaeptum de fine omnium praecaeptorum explicato  de libertate iustorum a lege de dile- ctione ex corde puro et fide non ficta et conscientia bona. SECVNDVM praecaedentes in te prophe- rias vt milites in illis bonam militiam.] In actibus apostolorum scribitur quod de timotheo multa bona  \pend\pstart teddebantur testimonia  quum paulus assumpsit eum secum. plus autem explicatur hic  dicendo que   \pend
\section*{CAPVT }
\section*{.II. }
\marginpar{[ p.141 ]}\pstart praeceflerunt prophetiae de timotheo. Er quantum ex his coniicitur, ab aliqua seu aliquibus spiritualibus personis praedicta fuerant multa bona spiritualia timotheo accessura, quae tam paulo quam timotheo nota- erant. et propterea animum timothei accendit ex talibus praecedentibus prophetiis. Er est ordo literae. vt fm praecedentes in te (seu super te seu ad te) prophetias milites in illis (prophetiis: hoc est in prophetaris per illas prophetias) bonam militiam  agas bonum militem. HABENS fidem et bonam conscientiam.] Applicat ad ipsum timotheum quae in genere dixerat de fide et conscientia bona, vt timotheus sit primus ibi habens fidem et conscientiam bonam: vtpote qui praesidebat ephesi. QVAM quidam repellentes. pro- qua quidam repulsa.] Et relatiusi qua refert bonam conscientiam: ad ostendendum quanti pendenda est bona conscientia  hoc est bona applicatio scientiae. CIRCA fidem naufragauerunt, pro. naufragium fece- tunt.] Ecce periculum ac damnum non se uatae bonae conscientiae inde siquidem iacturam incurrerunt fidei. Male siquidem applicando scientiam  incurritur error in fide : et sic paulatim perdit̃ fides. EX QVIBVS est hymenaeus et alexander] Nihil inuenio autentice scriptum de his duobus qui repulsa bona conscientia perdiderunt fidem. QVOS tradidi satanae.] corporaliter vexandos: vt de corinthio incestuoso expli- catur. VT discant non blasphemare.] Quemadmodum corinthium exconicando tradidit satanae sum car- nem vt spiritus saluus fieret, ita istos exconicando tradidit satanae vt discant huiusmodi poenis non blasphe- mate. Insinuat .n. hinc quod iesum christum blasphemabant isti perdita fide. Et forte isti erant doctores docen- tes aduersus christum postquam a fide ceciderunt. et propterea dicit vt discant non blasphemare.  \pend
\endnumbering\beginnumbering\section{CAPVT .II.}\pstart \huge\textbf{O}\normalsize BSECRO. pro. adhortor  igitur primum omnium.] Tractata necessitate quare reliquerat ti- motheum ephesi  incipit prosequi proposita. Vnde tanquam ad proposita exequenda rediens  illa- tiue dicit adhortor igitur. Proposuerat autem quod finis praecaepti est dilectio de corde puro etc.  et propterea ab executione dilectionis erga omnes inchoat. Primus antem et communissimus dile- ctionis actus erga omnes est oratio. et ideo tractans illam parriculam dilectio  hortatur ante omnia orari pro omnibus. FIERI obsecrationes. pro. preces.] petitorias. ORATIONES. ] hoc est mentis eleuationes in deum. POSTVLATIONES. pro. intercessiones. ] Quae differunt a precibus in hoc quod preces respiciunt rem petitam: intercessiones autem mediae sunt ad reconciliandum homines deo. GRATIARVM ACTIO- NES.] pro beneficiis accaeptis. PRO omnibus hominibus.] Ecce extensio dilectionis ad omnes ho- mines siue christianos siue non christianos. nam non dicit pro omnibus fratribus sed omnibus homini- bus. PRO regibus] qui tunc non erant christiani. Vult pro omnibus hominibus in genere  pro regi- bus autem in specie orari. Nec solum pro regibus sed ET OMNIBVS qui in sublimitate constituti sunt. ] Omnes temporales dominos compraehendit  generali nomine fublimitaris seu excellentiae. VT QVIETAM et tranquillam vitam agamus] Non concordant omnino interpretes circa haec vocabula. Verum parui re- fert tranquillam et quietam vitam an placidam et quietam seu tacitam dixeris. Id tamen aduertendum est quod quum dixisset tum pro omnibns homnibus tum pro regibus et dominis, prius more suo declarat postremum. s. quare pro regibus et in celsitudine constituris  vt quietam et tranquillam vitam etc.  Siquidem sermo pauli formalis est  quum dicit pro regibus et temporalibus dominis orandum: videlicet quatenus regunt et dominant̃ vt regimen eorum iustum sit. nam inde sequitur iste fructus vt tranquillam et quietam vitam agamus chri- stiani. IN OMNI pietate et castitate. pro. honestate.] Ne quietem et tranquillitatem vitae apperant chri- stiani absolute sed relatiue ad virtutum opera  adiungit virtutem pietatis erga parentes patriam etc.  et vniuersae morum grauiratis. Et dicit in omni, vt nulla pars pietatis  nulla pars honesti moris praetermittatur.  \pend\pstart HOC ENIM bonum est.] Pronomen hoc demonstrat orare pro omnibus hominibus. Et quod primo dixerat (.s. pro omnibus hominibus orandum) manifestat hoc esse bonum in se. ET accaeptum coram sal- uatore nostro deo. ] Non solum est in se bonum sed gratum deo. Quem nominat saluatorem, quia ex officio saluandi manifestat intentum  subiungens. QVI omnes homines vult saluos fieri et ad agnitionem veri- taris venire.] Non est hic speculatiue theologe sermo de voluntate bñplaciti : de qua scriptum est, omnia quae voluit dominus fecit in coelo et in terra et quod voluntati eius nullus porest resistere, vt clare patet ex eo quod vi- demus non omnes homines ad agnitionem veritatis venire. Sed est sermo de voluntate signi : qua deus propo- nit omnibus hominibus praecaepta salutis doctrinamque  euangelicant et illis imputatur qui nolunt venire ad veritaris cognirionem et nolunt salui fieri. de qua scriptum est, quotiens volui congregare filios tuos et no- luistis. Ex hoc enim quod deus salutem et agnitionem omnibus vult proponendo, monstratur gratum illi esse et pro omnibus oremus, et simul insinuatur quid est petendum orando pro omnibus  videlicet vt saluent̃. ⁋ VNVS enim deus. ] Manifestat quod deus vult omnes homines saluos fieri et ad agnitionem veritatis ve- nite, tum ex vnitate dei.ex hoc enim quod non nisi vnus est deus  omnium hominum curam illi vm incumbere manifestatur: et quum sit natura bonus, consequens est vt omnibus proponat salutem et agnitionem verita tis. Si namque  essent plures dii, cogitari posset quod vnus deus haberet curam saluandi aliquos homines  et alius haberet curam saluandi alios: sed vbi vnus tantum est deus, illi vni incumbit cura omnium. Tum ex vni- tate mediatoris  subiugendo. VNVS mediator dei et hominum homo christus iesus.] Si płutes essent mediatores inter deum et hominem, existimaretur quod vnus esset mediator pro quibusdam de alius pro aliis sed  \pend
\section*{PRIORIS }
\section*{AD TIMOTHEVM }\pstart ex quo vnus est mediator dei et hominum ad reconciliandum homines deo, illi vni incumbit mediare inter deum et omnes homines. Et propterea idem sequitur.s.quod deus vult omnes homines saluos fieri et ad agnitionem veritatis venite  ex hocipso quod vnum constituit mediatorem dei et hominum. Qui tanto ma- gis omnium hominum mediator est ad deum quanto ipse est homo  quanto ab ipso nihil humani alienum est. QVI dedit redemptionem semetipsum pro omnibus. ] Ab effectu manifestat quod etiam mediator vult omnes homines saluos fieri et ad agnitionem veritatis veniretex hocipso quod dedit semetipsum precium ad redimendum non aliquos sed omnes. Quod verissimum est non solum quantum ad sufficientiam sed etiam quantum ad efficientiam: pro quanto omnes redemit a morte corporis resuscitando omnes in nouissimo die. VCVIVS testimonium temporbus suis confirmatum est.] Superfluunt tres dictiones; cuius confirmatum est. Legendum siquidem est apposiriue. dedit redemptionem semetipsum pro omnibns appositiue testimo- nium temporibus suis. Vere testimonium diuinae voluntatis q vult omnes homines saluos fieri et ad agni- tionem veritatis venire  dedit christus dando seipsum redemptionem pro omnibus. dedit autem hoc testi- monium temporibus propnis definitis a diuina sapientia. Et addidit hoc, ne quisꝗ diceret tardum fuisse testimonium respectiue ad eos qui praecesserant aduentum christi omnis enim tam tarditas quam anticipatio excludirur dicendo tperibus suis. IN QVO. pro. in quod.) videlicet testimonium promulgandum mundo. POSITVS sum ego predicator pro. praeco. ] Officium preconis est publicare: et paulus publicabat testimo- nium datum a christo. ET apostolus.] Ne vsurpatum aut vile praeconis officium putaretur, adiunxit et apostolus  proculdubio christi. ¶ VERITATEM dico. ] Deficit. IN christomnon mentior.] Parenthe- sis est haec, explicans officium praeconis et apostoli paulum agere vtranque  complectendo partem : videlicet non mentiri et veritatem dicere christo teste. vt nulla falsitas mixta verirati intelligatur. DOCTOR gentium in fide et verirate.] Non solum praeco et apostolus sum veritatis, sed sum specialiter doctor gentium in fide non in pħia, sed non in fide vana  sed et veritate. Doceo gentes veritatem  non figmenta : verum non in eui- dentia scientifica sed in fide. et haec omnia ago ad publicandum testimonium quod dedit christus.  \pend\pstart \huge\textbf{V}\normalsize OLO ERGO viros. ] Tractat illam particulam  de corde puro. Continuat autem exercens officium doctoris gentium. ORARE in omni loco, leuantes puras manus sine ira et disceptatione.] Di- xerat dominus  si offers munus tuum ante altare et recordatus fueris quod frater tuus habet aliquod  aduersum te etc extendit paulus ad omnem locum idem praeceptum. Docet siquidem quod non solum ante altare sed in omni loco in quo ꝗs orauerit oportet leuari manus mundas sine ira et disceptatione. Nec est sermo de membris corporeis sed de operibus. impedimenta siquidem orationis vbicumque  oretur, impuriras operdĩ  ira et conten- tio sunt ¶ Et excluduntur ad literam duo crimina. Alterum auariciae occupantis, dicendo puras manus. communi siquidem vsu occupantes aliena dicimus non habere manus puras. Alterum discordiae cum proxi- mo: dicendo sine ira  intus, et disceptatione extra. haec enim duo purganda sunt ante orarionem, auaricia restituendo et discordia reconciliatione. Et meminit mundiciae ab his criminibns tantum: quia reliqua eri- mina interna poenitentia purgantur, haec autem exigunt etiam externam restitutionem et reconcilianone. ¶ SIMILITER et mulieres in habitu. pro amictu  ornato cum verecundia et sobrietate. pro. modestia, ornantes pro ornate  seipsas.] Quemadmodum a viris abstulit impedimenta orationis tanquam propria viris, ita a mulieribus tollit impedimenta orationis quae sunt velut propria mulieribus. Vnde quod dicit in amictu ornato, sic dicit quemadmodum dixerat in omni loco. Vtrumque  enim recensetur ad explicandum mo- dum seruandum  a viris quidem in quocumque  loco orauerint: a mulieribus vero in quocumque  amictu omnato. Et explicatur modus dicendo cum verecundia et modestia ornare seipsas. Vtraque  pars valde congruit mu- ueribus: ne ornatus sit inuerecundus, et similiter ne sit imoderatus. Vnde imoderationem ornatus explican- do subiungit. NON in tortis crinibus aut auro aut margartis.] Primiriuae ecclesiae tempori congrue- bat hoc praeceptum:eo quod tunc mulieres quae se christianas profitebantur  omnia huius modi a se abdica- bant tanquam superflua et vana. Vnde hodie quum multae principes et reginae quibus sum statum congruit or- nare se auro et margaritis  christianae sint, non contra pauli praeceptum his vtuntur. VEL veste preciosa pro. sumptuosa.] Dixerat in amictu omnato  et declarauerat imoderatum omatum:declarat modo ami- ctum, inhibendo vestem sumptuosam. Et quod temporale fuerit huiusmodi praeceptum  aurea vestis ceciliae te- statur. SED QVOD. pro. quae  decet mulieres promittentes. pro. profirentes  pietatem per opera bona.] Re- latiuum quae ad vestem refertur. et vult vestem cem talem quae deceat mulieres profirentes per bona opera (hoc est profitentes non tantum vetbis sed factis) pietatem:hoc est cultum diuinum. Dictio siquidem graeca non significat pietatem in genere, sed eam quae est erga deum. Perpende ex his omnibus impedimenta ora tionis in mulienbus vbicumque  orent esse contraria praedictorum: vtpote appetitum muliebrem deprimentia ad haec vana ne eleuetur illarum mens in deum. ¶ MVLIER in silentio discat cum omni subiectione.] Postquam cor purum tam in viris quam in mulieribus declarauit relatiue ad orationem, declarat in specie puritarem cor- dis in muliere relatine tum ad disciplinam et doctrinam  tum ad dominium. Et qm̃ mulier eget vt docearur, modum explicat tum quo ad linguam  dicendo in silentio. Non vult quod in ecelesia interroget doctorem do- centem, sed quod audiat quidem et discat silendo:si quid autem restat dubii  virum domi interroget. Tum quo ad reuetentiam cum omni subiectione  interna et externa ¶ DOCERE autem ] supple publice. MULIERE  \pend
\section*{CAPVT }
\section*{.III. }
\marginpar{[ p.142 ]}\pstart non permitto.] Et diximus publice  ad differentiam doctrinae domesticae : qua mater debet erudite filios et filias. NEQVE dominari pro. neque  autoritatem habere  in virum. ] Directe hoc respicit vxores : quas nefas est autoritarem habere in maritos. SED essein silentio.] vt non solum opere sed nec sermone auto- ritatem sibi vsurpet in virum. ADAM enim primus formatus est, deinde eua.] Exordine quo moyses narrat formationem adae et euae, probat virum praeefse mulieri et non econuerso. ET ADAM non est seductus, mulier aurem seducta in praeuaricatione fuit. ] Eadem vtitur historia : in qua mulier ipsa pro- fessa est se seductam a serpente. adam vero non se excusauit vt decaeptum  sed vt allectum ab vxore : vt ibi pa- tet. Praestat itaque  vir mulieri tum ordine formationis  tum vigore rationis. SALVABITVR autem per filiorum generationem si permanserit, pro. si manserint  in fide et dilectione et sanctificatione cum so- brietate pro. modestia.] Ordo literae est. saluabitur autem si manserint (ipsa et vir) per filiorum genera- tionem in fide etc. ¶ Meritorium salutis astruit esse coningium mulieri : ex hoc quod si ob filiorum genera- tionem manserit simul cum viro in fide etc. saluabitur ipsa mulier. Et per hoc intellige explicari pro- prium munus sexus. Non docet paulus salutem mulieris pendere a generatione filiorum (quum facilius saluetur virgo quam coningata) sed docet quod vtendo officio sexus  habet vnde ex ipso vsu sexus saluetur : hoc est quod ad salutem eis operatur si ob filiorum generationem manserint cum viris suis in fide etc.  Et si vis oemm scrupulum ab hac lectione excludere, ordina saluabitur pro predicato totius enunciationis : ordinando li- teram sic.si permanserint autem per filiorum generationem in fide et dilectione et sanctificatione cum mode- stia  saluabitur. Sic .n. clare significatur quod non pendet salus mulieris a filiorum generatione, sed quod acquirit̃ ei salus si sic manserint.¶ Et nihilominus aduerte quod non de muliere sed de vxore est sermo: de qua dixerat neque  autoritatem habere in virum. de ipsa .n. subiungit saluabitur asit etc.  ¶ Et meminit specialiter sancti- taris ac modestiae : quoniam tam mundicia in vsu coningii quam moderatio in vsu rerum et gubernatione familiae  peculiaris esse debet coniugatis.  \pend
\endnumbering\beginnumbering\section{CAPVT .III.}\pstart \huge\textbf{F}\normalsize IDELIS sermo. ] Prosequitur paulus particulam de corde puro  quo ad constitutos in ec- clesiasticis ministeriis. Praeponit aurem huic tractarui fidelis sermo  ad significandam veri- tatem dicendorum : ne quis imputaret paulo quod nimis exigeret ex proprio capite in episcopis et diaconis. Totum hoc excluditur  dicendo fidelis sermo. significatur enim quod est fermo fide- lis domino, quod est iuxta fidem quam debet paulus domino  et non sum propriam sententiam. SI QVIS episcopatum desiderat  bonum opus desiderat.] Initium fidelis sermonis est si quis superintendentiam cupit, bonum (seu praeclarum) opus desiderat. Non dicit si quis dignitatem  si quis gradum  si quis pro- uentus  si quis honores  si quis gloriam episcopatus : sed si quis ipsum episcopatum (hoc est ipsam superin tendentiam  ipsum superintendere) desiderat, proculdubio desiderat opus quum superintendere sit opus, non qualecunque  sed opus bonum ex suo genere  immo opus praeclarum ex proprio genere. Nec hoc eget probatione aliqua : quoniam manifeste constat sic esse. Nullum itaque  peccatum est ex proprio genere  desi- derare episcopatum: hoc est ipsum superintendere. Secus autem est de appetitu gradus  honoris  prouen- tuum et huiusmodi. nihil enim horum est opus bonum: quum nec sit opus. OPORTET enim. pro. igitur  episcopum irrepraehensibilem esse.] Ex ratione excellentis bonitatis huiusmodi operis quod est supin renderejillariue dicit oporter igirur episcopum esse irrepraehensibilem. Et haec est prima conditio episcopi relatiue ad malos mores  vt nihil in eo sit repraehensibile. Et vere ex eo quod aliis superintendit repraehenden- dis, manifeste infertur quod oportet ipsum esse irrepraehensibilem. qua enim fronte porest alios repraehendere qui repraehensibilis est. diceretur namque  ei  medice cura teipsum. VNIVS vxoris virum.] Secunda condi- tio negatiue intelligenda est : hoc est non plurium vxorum virum. Non enim exigit paulus ab episcopo quod fuerit aut sit coniugatus (qui consulit meliorem esse coelibatum  scribendo ad corin.) sed exigit quod si habet vrorem sit vir vnius vxoris. Sic .n. litera sonat  dicendo esse. Quo contra de vidua inferius dicet in ca. 5. non praesenti sed preterito vtens tempore  quae fuerit vnius viri vxor. Et dixit hoc, quia tunc temporis multi plures vxores habebant  imitantes patres veteris testamenti : quum nullibi legamus hoc prohibitum. SOBRIVM pro. vigilem.P Tertia conditio vigilantiam summe necessariam supintendenti explicat. PRVDENTEM pro. modestum.] Quarta conditio summe quoque  est necessaria superintendenti. vt vigilantia comitem habeat moderationem : vt ipse sic sit vigil vt tñ sit modestus  et sic modestus vt tñ sit vigil. ORNATVM.] proculdubio compositione gestuum  actionum  eloquii  et aliis bonis moribus. PVDICVM.] Superfluit pudicum. Sexto .n. loco dr̃. HOSPITALEM.P Rara est haec hodie virtus. DOCTOREM.] non laurea, sed scientia sufficiente ad docendum. NON vinolentum. pro. non vinosum.] hoc est deditum vino. NON percussorem.] Quamuis paulus de iniusta pcussione loquatur, turpe tñ est epreo eriam iuste percutere, debet .n. aliena non propria manu percutere delinquentes. ¶ Post non percussorem  deficit vna particula. s. NON tut- pilucrem. ] Haec est .n. decima conditio in textu pauli : clare sonans quod epres alienus esse a turpi lucro debet. Quod si seruaret̃, non audirent̃ lucra tam turpia vt turpe sit ea dicere. Nec loquor de simonia, sed proprie de turpi lucro sparso nimis in epreis ecclesiarum christi. SED modestum. pro. sed mitem.] Esse mitem plus est quam esse mansuetum mitis enim non solum iram temperat, sed dulcedinem quandam adhibet. NON litigo-  \pend
\endnumbering
\end{pages}
\end{document}
        